% Options for packages loaded elsewhere
\PassOptionsToPackage{unicode}{hyperref}
\PassOptionsToPackage{hyphens}{url}
\documentclass[
  french,
]{article}
\usepackage{xcolor}
\usepackage[margin=2cm]{geometry}
\usepackage{amsmath,amssymb}
\setcounter{secnumdepth}{-\maxdimen} % remove section numbering
\usepackage{iftex}
\ifPDFTeX
  \usepackage[T1]{fontenc}
  \usepackage[utf8]{inputenc}
  \usepackage{textcomp} % provide euro and other symbols
\else % if luatex or xetex
  \usepackage{unicode-math} % this also loads fontspec
  \defaultfontfeatures{Scale=MatchLowercase}
  \defaultfontfeatures[\rmfamily]{Ligatures=TeX,Scale=1}
\fi
\usepackage{lmodern}
\ifPDFTeX\else
  % xetex/luatex font selection
\fi
% Use upquote if available, for straight quotes in verbatim environments
\IfFileExists{upquote.sty}{\usepackage{upquote}}{}
\IfFileExists{microtype.sty}{% use microtype if available
  \usepackage[]{microtype}
  \UseMicrotypeSet[protrusion]{basicmath} % disable protrusion for tt fonts
}{}
\makeatletter
\@ifundefined{KOMAClassName}{% if non-KOMA class
  \IfFileExists{parskip.sty}{%
    \usepackage{parskip}
  }{% else
    \setlength{\parindent}{0pt}
    \setlength{\parskip}{6pt plus 2pt minus 1pt}}
}{% if KOMA class
  \KOMAoptions{parskip=half}}
\makeatother
\usepackage{color}
\usepackage{fancyvrb}
\newcommand{\VerbBar}{|}
\newcommand{\VERB}{\Verb[commandchars=\\\{\}]}
\DefineVerbatimEnvironment{Highlighting}{Verbatim}{commandchars=\\\{\}}
% Add ',fontsize=\small' for more characters per line
\newenvironment{Shaded}{}{}
\newcommand{\AlertTok}[1]{\textcolor[rgb]{1.00,0.00,0.00}{\textbf{#1}}}
\newcommand{\AnnotationTok}[1]{\textcolor[rgb]{0.38,0.63,0.69}{\textbf{\textit{#1}}}}
\newcommand{\AttributeTok}[1]{\textcolor[rgb]{0.49,0.56,0.16}{#1}}
\newcommand{\BaseNTok}[1]{\textcolor[rgb]{0.25,0.63,0.44}{#1}}
\newcommand{\BuiltInTok}[1]{\textcolor[rgb]{0.00,0.50,0.00}{#1}}
\newcommand{\CharTok}[1]{\textcolor[rgb]{0.25,0.44,0.63}{#1}}
\newcommand{\CommentTok}[1]{\textcolor[rgb]{0.38,0.63,0.69}{\textit{#1}}}
\newcommand{\CommentVarTok}[1]{\textcolor[rgb]{0.38,0.63,0.69}{\textbf{\textit{#1}}}}
\newcommand{\ConstantTok}[1]{\textcolor[rgb]{0.53,0.00,0.00}{#1}}
\newcommand{\ControlFlowTok}[1]{\textcolor[rgb]{0.00,0.44,0.13}{\textbf{#1}}}
\newcommand{\DataTypeTok}[1]{\textcolor[rgb]{0.56,0.13,0.00}{#1}}
\newcommand{\DecValTok}[1]{\textcolor[rgb]{0.25,0.63,0.44}{#1}}
\newcommand{\DocumentationTok}[1]{\textcolor[rgb]{0.73,0.13,0.13}{\textit{#1}}}
\newcommand{\ErrorTok}[1]{\textcolor[rgb]{1.00,0.00,0.00}{\textbf{#1}}}
\newcommand{\ExtensionTok}[1]{#1}
\newcommand{\FloatTok}[1]{\textcolor[rgb]{0.25,0.63,0.44}{#1}}
\newcommand{\FunctionTok}[1]{\textcolor[rgb]{0.02,0.16,0.49}{#1}}
\newcommand{\ImportTok}[1]{\textcolor[rgb]{0.00,0.50,0.00}{\textbf{#1}}}
\newcommand{\InformationTok}[1]{\textcolor[rgb]{0.38,0.63,0.69}{\textbf{\textit{#1}}}}
\newcommand{\KeywordTok}[1]{\textcolor[rgb]{0.00,0.44,0.13}{\textbf{#1}}}
\newcommand{\NormalTok}[1]{#1}
\newcommand{\OperatorTok}[1]{\textcolor[rgb]{0.40,0.40,0.40}{#1}}
\newcommand{\OtherTok}[1]{\textcolor[rgb]{0.00,0.44,0.13}{#1}}
\newcommand{\PreprocessorTok}[1]{\textcolor[rgb]{0.74,0.48,0.00}{#1}}
\newcommand{\RegionMarkerTok}[1]{#1}
\newcommand{\SpecialCharTok}[1]{\textcolor[rgb]{0.25,0.44,0.63}{#1}}
\newcommand{\SpecialStringTok}[1]{\textcolor[rgb]{0.73,0.40,0.53}{#1}}
\newcommand{\StringTok}[1]{\textcolor[rgb]{0.25,0.44,0.63}{#1}}
\newcommand{\VariableTok}[1]{\textcolor[rgb]{0.10,0.09,0.49}{#1}}
\newcommand{\VerbatimStringTok}[1]{\textcolor[rgb]{0.25,0.44,0.63}{#1}}
\newcommand{\WarningTok}[1]{\textcolor[rgb]{0.38,0.63,0.69}{\textbf{\textit{#1}}}}
\usepackage{longtable,booktabs,array}
\usepackage{caption}
\captionsetup[table]{skip=6pt}
\newcounter{none} % for unnumbered tables
\usepackage{calc} % for calculating minipage widths
% Correct order of tables after \paragraph or \subparagraph
\usepackage{etoolbox}
\makeatletter
\patchcmd\longtable{\par}{\if@noskipsec\mbox{}\fi\par}{}{}
\makeatother
\ifLuaTeX
\usepackage[bidi=basic,shorthands=off]{babel}
\else
\usepackage[bidi=default,shorthands=off]{babel}
\fi
\ifLuaTeX
  \usepackage{selnolig} % disable illegal ligatures
\fi
\setlength{\emergencystretch}{3em} % prevent overfull lines
\providecommand{\tightlist}{%
  \setlength{\itemsep}{0pt}\setlength{\parskip}{0pt}}
\usepackage{bookmark}
\IfFileExists{xurl.sty}{\usepackage{xurl}}{} % add URL line breaks if available
\urlstyle{same}
% fallback for those not using the hyperref driver hyperxmp:
\makeatletter
\@ifundefined{xmpquote}{\newcommand{\xmpquote}[1]{#1}}{}
\makeatother
\hypersetup{
  pdflang={fr},
  hidelinks,
  pdfcreator={LaTeX via pandoc}}

\author{}
\date{}

\begin{document}

\section{Tutoriel Pratique Exhaustif COVOF 2D (Fortran 2018 /
OpenMPI)}\label{tutoriel-pratique-exhaustif-covof-2d-fortran-2018-openmpi}

Ce document constitue le \textbf{guide d'apprentissage complet pas à
pas} pour configurer, compiler, exécuter, analyser et développer des cas
de calcul en mécanique des fluides numérique avec la version
\textbf{Modern Fortran 2018 modulaire} du solveur \textbf{COVOF 2D}.

\begin{center}\rule{0.5\linewidth}{0.5pt}\end{center}

\subsection{Table des Matières}\label{table-des-matiuxe8res}

\begin{enumerate}
\def\labelenumi{\arabic{enumi}.}
\tightlist
\item
  \hyperref[1-introduction-et-notions-fondamentales]{Introduction et
  Notions Fondamentales}
\item
  \hyperref[2-installation-et-environnement-de-compilation]{Installation
  et Environnement de Compilation}
\item
  \hyperref[3-structure-et-syntaxe-dun-cas-test-casesjson]{Structure et
  Syntaxe d'un Cas Test (\texttt{cases/*.json})}
\item
  \hyperref[4-tutoriel-1--uxe9coulement-de-poiseuille-et-transitoire-de-szymanski]{Tutoriel
  1 : Écoulement de Poiseuille et Transitoire de Szymanski}
\item
  \hyperref[5-tutoriel-2--cavituxe9-entrauxeenuxe9e-uxe0-reynolds-re--1000]{Tutoriel
  2 : Cavité Entraînée à Reynolds Re = 1000}
\item
  \hyperref[6-tutoriel-3--uxe9coulement-sur-marche-descendante-armaly-re--100]{Tutoriel
  3 : Écoulement sur Marche Descendante (Armaly Re = 100)}
\item
  \hyperref[7-tutoriel-4--convection-naturelle-thermique-de-vahl-davis-ra--10ux2074]{Tutoriel
  4 : Convection Naturelle Thermique de Vahl Davis (Ra = 10⁴)}
\item
  \hyperref[8-tutoriel-5--uxe9coulement-diphasique-vof--rupture-de-barrage]{Tutoriel
  5 : Écoulement Diphasique VOF --- Rupture de Barrage}
\item
  \hyperref[9-tutoriel-6--tension-de-surface-et-saut-de-pression-de-laplace]{Tutoriel
  6 : Tension de Surface et Saut de Pression de Laplace}
\item
  \hyperref[10-tutoriel-7--uxe9coulement-autour-dun-cylindre-par-puxe9nalisation-de-brinkman]{Tutoriel
  7 : Écoulement Autour d'un Cylindre par Pénalisation de Brinkman}
\item
  \hyperref[11-post-traitement-avancuxe9-avec-paraview]{Post-Traitement
  Avancé avec ParaView}
\item
  \hyperref[12-automatisation-python-et-uxe9tudes-de-convergence]{Automatisation
  Python et Études de Convergence}
\item
  \hyperref[13-cruxe9er-un-nouveau-cas-test-uxe0-partir-de-zuxe9ro]{Créer
  un Nouveau Cas Test à Partir de Zéro}
\end{enumerate}

\begin{center}\rule{0.5\linewidth}{0.5pt}\end{center}

\subsection{1. Introduction et Notions
Fondamentales}\label{introduction-et-notions-fondamentales}

\subsubsection{1.1. Maillage Décalé MAC et Convention Mémoire
Fortran}\label{maillage-duxe9caluxe9-mac-et-convention-muxe9moire-fortran}

Le solveur COVOF 2D repose sur la discrétisation spatiale sur grille
décalée de Harlow \& Welch (1965) : * \textbf{Composante \(u\) de
vitesse} : discrétisée aux faces verticales \((i \pm \frac{1}{2}, j)\).
* \textbf{Composante \(v\) de vitesse} : discrétisée aux faces
horizontales \((i, j \pm \frac{1}{2})\). * \textbf{Champs scalaires
(\(p, T, C\))} : positionnés au centre des cellules \((i, j)\).

\begin{verbatim}
                  v(i, j+1/2)
            +----------▲----------+
            |          |          |
            |          |          |
u(i-1/2, j) ►     p(i, j)         ► u(i+1/2, j)
            |     T(i, j)         |
            |     C(i, j)         |
            |          |          |
            +----------▲----------+
                  v(i, j-1/2)
\end{verbatim}

En \textbf{Fortran 2018}, les tableaux 2D sont stockés en \textbf{ordre
colonne-major} (\emph{column-major order}). Pour garantir des accès
mémoire contigus et un alignement optimal pour la vectorisation SIMD
(AVX2/AVX-512), toutes les boucles du code respectent scrupuleusement la
hiérarchie :

\begin{Shaded}
\begin{Highlighting}[]
\KeywordTok{do}\NormalTok{ j }\KeywordTok{=} \DecValTok{1}\NormalTok{, ny}
   \KeywordTok{do}\NormalTok{ i }\KeywordTok{=} \DecValTok{1}\NormalTok{, nx}
      \CommentTok{! Accès contigu en mémoire à array(i, j)}
\NormalTok{      res(i, j) }\KeywordTok{=}\NormalTok{ array(i}\KeywordTok{+}\DecValTok{1}\NormalTok{, j) }\KeywordTok{{-}} \FloatTok{2.0\_dp} \KeywordTok{*}\NormalTok{ array(i, j) }\KeywordTok{+}\NormalTok{ array(i}\KeywordTok{{-}}\DecValTok{1}\NormalTok{, j)}
   \KeywordTok{end do}
\KeywordTok{end do}
\end{Highlighting}
\end{Shaded}

\subsubsection{1.2. Formulation par Lagrangien Augmenté et Itérations
d'Uzawa}\label{formulation-par-lagrangien-augmentuxe9-et-ituxe9rations-duzawa}

Pour résoudre les équations de Navier-Stokes incompressibles avec une
précision machine sur la divergence \(\nabla \cdot \mathbf{u} = 0\),
COVOF 2D utilise la méthode du Lagrangien Augmenté :
\[\rho \left( \frac{\partial \mathbf{u}}{\partial t} + (\mathbf{u} \cdot \nabla) \mathbf{u} \right) - \nabla \cdot (\mu (\nabla \mathbf{u} + \nabla^T \mathbf{u})) - r_{\text{al}} \nabla(\nabla \cdot \mathbf{u}) + \nabla p = \mathbf{f}\]
suivie de l'actualisation de pression d'Uzawa :
\[p^{m+1} = p^m - r_{\text{al}} (\nabla \cdot \mathbf{u}^{m+1})\]

\subsubsection{1.3. Politique
Zero-Allocation}\label{politique-zero-allocation}

Toutes les structures de données allouables (\texttt{flow\_fields\_t},
matrices implicites temporaires dans \texttt{navier\_stokes\_t}) sont
dimensionnées une seule fois lors de la phase d'initialisation
(\texttt{simulation\_init}). La boucle en temps s'exécute avec
\textbf{zéro appel aux allocateurs système
(\texttt{allocate}/\texttt{deallocate})}, ce qui élimine les temps morts
et garantit un déterminisme temporel parfait.

\begin{center}\rule{0.5\linewidth}{0.5pt}\end{center}

\subsection{2. Installation et Environnement de
Compilation}\label{installation-et-environnement-de-compilation}

\subsubsection{2.1. Vérification de la Chaîne
d'Outils}\label{vuxe9rification-de-la-chauxeene-doutils}

Vérifiez la présence du compilateur MPI Fortran et de CMake :

\begin{Shaded}
\begin{Highlighting}[]
\ExtensionTok{mpifort} \AttributeTok{{-}{-}version}
\FunctionTok{cmake} \AttributeTok{{-}{-}version}
\end{Highlighting}
\end{Shaded}

Le code est strictement conforme au standard \textbf{Fortran 2018} et
compatible avec GCC (\texttt{gfortran} 9+) et Intel oneAPI (\texttt{ifx}
/ \texttt{ifort}).

\subsubsection{2.2. Méthode A : Compilation avec CMake
(Recommandé)}\label{muxe9thode-a-compilation-avec-cmake-recommanduxe9}

\begin{Shaded}
\begin{Highlighting}[]
\BuiltInTok{cd}\NormalTok{ COVOF\_2D{-}FORTRAN\_fortran}
\FunctionTok{cmake} \AttributeTok{{-}B}\NormalTok{ build }\AttributeTok{{-}DCMAKE\_BUILD\_TYPE}\OperatorTok{=}\NormalTok{Release}
\FunctionTok{cmake} \AttributeTok{{-}{-}build}\NormalTok{ build }\AttributeTok{{-}j4}
\end{Highlighting}
\end{Shaded}

L'exécutable est généré directement dans \texttt{bin/covof2d\_fortran}.

\subsubsection{\texorpdfstring{2.3. Méthode B : Compilation Directe avec
\texttt{mpifort}}{2.3. Méthode B : Compilation Directe avec mpifort}}\label{muxe9thode-b-compilation-directe-avec-mpifort}

Si CMake n'est pas installé sur votre cluster, compilez directement les
modules ordonnés :

\begin{Shaded}
\begin{Highlighting}[]
\ExtensionTok{mpifort} \AttributeTok{{-}O3} \AttributeTok{{-}march}\OperatorTok{=}\NormalTok{native }\DataTypeTok{\textbackslash{}}
\NormalTok{  src/module\_precision.f90 }\DataTypeTok{\textbackslash{}}
\NormalTok{  src/module\_constants.f90 }\DataTypeTok{\textbackslash{}}
\NormalTok{  src/module\_json.f90 }\DataTypeTok{\textbackslash{}}
\NormalTok{  src/module\_config.f90 }\DataTypeTok{\textbackslash{}}
\NormalTok{  src/module\_mesh.f90 }\DataTypeTok{\textbackslash{}}
\NormalTok{  src/module\_fields.f90 }\DataTypeTok{\textbackslash{}}
\NormalTok{  src/module\_mpi\_comm.f90 }\DataTypeTok{\textbackslash{}}
\NormalTok{  src/module\_interface\_vof.f90 }\DataTypeTok{\textbackslash{}}
\NormalTok{  src/module\_navier\_stokes.f90 }\DataTypeTok{\textbackslash{}}
\NormalTok{  src/module\_linear\_solvers.f90 }\DataTypeTok{\textbackslash{}}
\NormalTok{  src/module\_vtk\_writer.f90 }\DataTypeTok{\textbackslash{}}
\NormalTok{  src/module\_simulation.f90 }\DataTypeTok{\textbackslash{}}
\NormalTok{  src/covof2d\_main.f90 }\DataTypeTok{\textbackslash{}}
  \AttributeTok{{-}o}\NormalTok{ bin/covof2d\_fortran}
\end{Highlighting}
\end{Shaded}

\begin{center}\rule{0.5\linewidth}{0.5pt}\end{center}

\subsection{\texorpdfstring{3. Structure et Syntaxe d'un Cas Test
(\texttt{cases/*.json})}{3. Structure et Syntaxe d'un Cas Test (cases/*.json)}}\label{structure-et-syntaxe-dun-cas-test-cases.json}

Le module \texttt{module\_json.f90} intègre un analyseur lexical et
syntaxique JSON natif en Fortran. Chaque cas de calcul est paramétré par
un fichier \texttt{.json} structuré en 7 sections :

\begin{Shaded}
\begin{Highlighting}[]
\FunctionTok{\{}
  \DataTypeTok{"mode"}\FunctionTok{:} \StringTok{"NAVIER\_STOKES"}\FunctionTok{,}
  \DataTypeTok{"mesh"}\FunctionTok{:} \FunctionTok{\{}
    \DataTypeTok{"nx"}\FunctionTok{:} \DecValTok{64}\FunctionTok{,}
    \DataTypeTok{"ny"}\FunctionTok{:} \DecValTok{64}\FunctionTok{,}
    \DataTypeTok{"lx"}\FunctionTok{:} \FloatTok{1.0}\FunctionTok{,}
    \DataTypeTok{"ly"}\FunctionTok{:} \FloatTok{1.0}
  \FunctionTok{\},}
  \DataTypeTok{"physics"}\FunctionTok{:} \FunctionTok{\{}
    \DataTypeTok{"enable\_two\_phase"}\FunctionTok{:} \KeywordTok{false}\FunctionTok{,}
    \DataTypeTok{"enable\_convection"}\FunctionTok{:} \KeywordTok{true}\FunctionTok{,}
    \DataTypeTok{"enable\_energy"}\FunctionTok{:} \KeywordTok{false}\FunctionTok{,}
    \DataTypeTok{"gx"}\FunctionTok{:} \FloatTok{0.0}\FunctionTok{,}
    \DataTypeTok{"gy"}\FunctionTok{:} \FloatTok{0.0}\FunctionTok{,}
    \DataTypeTok{"fluid1"}\FunctionTok{:} \FunctionTok{\{} \DataTypeTok{"rho"}\FunctionTok{:} \FloatTok{1.0}\FunctionTok{,} \DataTypeTok{"mu"}\FunctionTok{:} \FloatTok{0.001} \FunctionTok{\}}
  \FunctionTok{\},}
  \DataTypeTok{"time"}\FunctionTok{:} \FunctionTok{\{}
    \DataTypeTok{"dt"}\FunctionTok{:} \FloatTok{0.005}\FunctionTok{,}
    \DataTypeTok{"t\_final"}\FunctionTok{:} \FloatTok{5.0}\FunctionTok{,}
    \DataTypeTok{"cfl\_max"}\FunctionTok{:} \FloatTok{0.5}\FunctionTok{,}
    \DataTypeTok{"adaptive\_dt"}\FunctionTok{:} \KeywordTok{false}
  \FunctionTok{\},}
  \DataTypeTok{"boundary\_conditions"}\FunctionTok{:} \FunctionTok{\{}
    \DataTypeTok{"left"}\FunctionTok{:}   \FunctionTok{\{} \DataTypeTok{"type"}\FunctionTok{:} \StringTok{"DIRICHLET"}\FunctionTok{,} \DataTypeTok{"u"}\FunctionTok{:} \FloatTok{0.0}\FunctionTok{,} \DataTypeTok{"v"}\FunctionTok{:} \FloatTok{0.0} \FunctionTok{\},}
    \DataTypeTok{"right"}\FunctionTok{:}  \FunctionTok{\{} \DataTypeTok{"type"}\FunctionTok{:} \StringTok{"DIRICHLET"}\FunctionTok{,} \DataTypeTok{"u"}\FunctionTok{:} \FloatTok{0.0}\FunctionTok{,} \DataTypeTok{"v"}\FunctionTok{:} \FloatTok{0.0} \FunctionTok{\},}
    \DataTypeTok{"bottom"}\FunctionTok{:} \FunctionTok{\{} \DataTypeTok{"type"}\FunctionTok{:} \StringTok{"DIRICHLET"}\FunctionTok{,} \DataTypeTok{"u"}\FunctionTok{:} \FloatTok{0.0}\FunctionTok{,} \DataTypeTok{"v"}\FunctionTok{:} \FloatTok{0.0} \FunctionTok{\},}
    \DataTypeTok{"top"}\FunctionTok{:}    \FunctionTok{\{} \DataTypeTok{"type"}\FunctionTok{:} \StringTok{"DIRICHLET"}\FunctionTok{,} \DataTypeTok{"u"}\FunctionTok{:} \FloatTok{1.0}\FunctionTok{,} \DataTypeTok{"v"}\FunctionTok{:} \FloatTok{0.0} \FunctionTok{\}}
  \FunctionTok{\},}
  \DataTypeTok{"solver"}\FunctionTok{:} \FunctionTok{\{}
    \DataTypeTok{"type"}\FunctionTok{:} \StringTok{"AUGMENTED\_LAGRANGIAN"}\FunctionTok{,}
    \DataTypeTok{"tol"}\FunctionTok{:} \FloatTok{1.0e{-}7}\FunctionTok{,}
    \DataTypeTok{"max\_iter"}\FunctionTok{:} \DecValTok{500}\FunctionTok{,}
    \DataTypeTok{"r\_al"}\FunctionTok{:} \FloatTok{100.0}\FunctionTok{,}
    \DataTypeTok{"uzawa\_iter"}\FunctionTok{:} \DecValTok{1}\FunctionTok{,}
    \DataTypeTok{"linear\_solver"}\FunctionTok{:} \StringTok{"BICGSTAB2"}\FunctionTok{,}
    \DataTypeTok{"use\_precond"}\FunctionTok{:} \KeywordTok{true}\FunctionTok{,}
    \DataTypeTok{"precond\_type"}\FunctionTok{:} \StringTok{"JACOBI"}
  \FunctionTok{\},}
  \DataTypeTok{"output"}\FunctionTok{:} \FunctionTok{\{}
    \DataTypeTok{"output\_dir"}\FunctionTok{:} \StringTok{"output\_cavity\_re1000"}\FunctionTok{,}
    \DataTypeTok{"case\_name"}\FunctionTok{:} \StringTok{"cavity\_re1000"}\FunctionTok{,}
    \DataTypeTok{"vtk\_interval"}\FunctionTok{:} \DecValTok{20}\FunctionTok{,}
    \DataTypeTok{"log\_interval"}\FunctionTok{:} \DecValTok{10}\FunctionTok{,}
    \DataTypeTok{"write\_csv\_history"}\FunctionTok{:} \KeywordTok{true}
  \FunctionTok{\}}
\FunctionTok{\}}
\end{Highlighting}
\end{Shaded}

\subsubsection{Table des Paramètres
Principaux}\label{table-des-paramuxe8tres-principaux}

{\def\LTcaptype{none} % do not increment counter
\begin{longtable}[]{@{}
  >{\raggedright\arraybackslash}p{(\linewidth - 6\tabcolsep) * \real{0.2500}}
  >{\raggedright\arraybackslash}p{(\linewidth - 6\tabcolsep) * \real{0.2500}}
  >{\raggedright\arraybackslash}p{(\linewidth - 6\tabcolsep) * \real{0.2500}}
  >{\raggedright\arraybackslash}p{(\linewidth - 6\tabcolsep) * \real{0.2500}}@{}}
\toprule\noalign{}
\begin{minipage}[b]{\linewidth}\raggedright
Section
\end{minipage} & \begin{minipage}[b]{\linewidth}\raggedright
Clé
\end{minipage} & \begin{minipage}[b]{\linewidth}\raggedright
Valeurs / Formats
\end{minipage} & \begin{minipage}[b]{\linewidth}\raggedright
Signification Physique
\end{minipage} \\
\midrule\noalign{}
\endhead
\bottomrule\noalign{}
\endlastfoot
\texttt{mode} & \texttt{mode} & \texttt{"NAVIER\_STOKES"},
\texttt{"HEAT"} & Physique principale résolue \\
\texttt{mesh} & \texttt{nx}, \texttt{ny} & Entiers (\(\ge 16\)) & Nombre
de cellules le long des axes \(x\) et \(y\) \\
\texttt{mesh} & \texttt{lx}, \texttt{ly} & Réels en mètres & Dimensions
géométriques totales \\
\texttt{physics} & \texttt{enable\_two\_phase} & \texttt{true},
\texttt{false} & Active le suivi d'interface VOF \\
\texttt{physics} & \texttt{sigma} & Réel en \(\text{N/m}\) & Tension
interfaciale (modèle CSF) \\
\texttt{boundary} & \texttt{type} & \texttt{"DIRICHLET"},
\texttt{"NEUMANN"}, \texttt{"PERIODIC"}, \texttt{"FREESLIP"} & Condition
à la limite sur chaque frontière \\
\texttt{solver} & \texttt{linear\_solver} & \texttt{"BICGSTAB2"},
\texttt{"BICGSTAB"}, \texttt{"CG"}, \texttt{"GMRES"} & Algorithme
itératif Matrix-Free \\
\texttt{solver} & \texttt{r\_al} & Réel (\(\ge 0\)) & Coefficient de
pénalité de divergence \\
\end{longtable}
}

\begin{center}\rule{0.5\linewidth}{0.5pt}\end{center}

\subsection{4. Tutoriel 1 : Écoulement de Poiseuille et Transitoire de
Szymanski}\label{tutoriel-1-uxe9coulement-de-poiseuille-et-transitoire-de-szymanski}

Ce cas valide l'établissement transitoire d'un profil de vitesse sous
l'application brutale d'une force volumique motrice axiale
\(g_x = 0.0555\text{ m/s}^2\) dans un canal plan de demi-largeur
\(e = 0.08\text{ m}\).

\subsubsection{\texorpdfstring{4.1. Configuration
(\texttt{cases/poiseuille\_unsteady.json})}{4.1. Configuration (cases/poiseuille\_unsteady.json)}}\label{configuration-casespoiseuille_unsteady.json}

\begin{itemize}
\tightlist
\item
  Propriétés du fluide : \(\rho = 1.0\text{ kg/m}^3\),
  \(\mu = 0.01\text{ Pa}\cdot\text{s}\).
\item
  Conditions aux limites : Périodique en \(x\), adhérence (\(u=0, v=0\))
  en haut et en bas.
\item
  Pénalisation \(r_{\text{al}} = 0.0\) (car la vitesse est purement
  solénoïdale et \(p \equiv 0\)).
\end{itemize}

\subsubsection{4.2. Exécution
Parallèle}\label{exuxe9cution-paralluxe8le}

\begin{Shaded}
\begin{Highlighting}[]
\ExtensionTok{mpirun} \AttributeTok{{-}np}\NormalTok{ 2 bin/covof2d\_fortran cases/poiseuille\_unsteady.json}
\end{Highlighting}
\end{Shaded}

\subsubsection{4.3. Solution Analytique (Série de
Fourier)}\label{solution-analytique-suxe9rie-de-fourier}

La vitesse axiale exacte s'écrit sous forme de série infinie :
\[u(z, t) = \frac{G e^2}{2\mu} \left( 1 - \frac{z^2}{e^2} \right) + \sum_{n=0}^{\infty} a_n \exp\left( - \frac{(2n+1)^2 \pi^2 \mu}{4 \rho e^2} t \right) \cos\left( \frac{(2n+1) \pi z}{2e} \right)\]
avec \(a_n = -\frac{16 G e^2 (-1)^n}{\mu \pi^3 (2n+1)^3}\) et
\(z = y - e\).

À \(t = 1.0\text{ s}\), la solution numérique Fortran concorde avec la
solution analytique avec une erreur
\(L_1 = 1.49 \times 10^{-5}\text{ m/s}\) (très inférieure à la cible
Référence Académique de \(3.13 \times 10^{-4}\text{ m/s}\)).

\begin{center}\rule{0.5\linewidth}{0.5pt}\end{center}

\subsection{5. Tutoriel 2 : Cavité Entraînée à Reynolds Re =
1000}\label{tutoriel-2-cavituxe9-entrauxeenuxe9e-uxe0-reynolds-re-1000}

Le benchmark canonique de Ghia et al.~(1982) valide la discrétisation
des termes convectifs et visqueux à haut nombre de Reynolds.

\begin{verbatim}
       u_top = 1.0 m/s, v_top = 0.0 m/s
    +--------------------------------+
    |                                |
    |           VORTEX               |
    |          PRIMAIRE              |
    |                                |
    |  Vortex                  Vortex|
    |  coin bas-gauche   coin bas-dr.|
    +--------------------------------+
             Parois fixes (u=0, v=0)
\end{verbatim}

\subsubsection{\texorpdfstring{5.1. Configuration
(\texttt{cases/cavity\_re1000.json})}{5.1. Configuration (cases/cavity\_re1000.json)}}\label{configuration-casescavity_re1000.json}

\begin{itemize}
\tightlist
\item
  Domaine unitaire \(1.0\text{ m} \times 1.0\text{ m}\), maillage fin
  \(128 \times 128\).
\item
  Fluide : \(\rho = 1.0\text{ kg/m}^3\),
  \(\mu = 0.001\text{ Pa}\cdot\text{s} \implies Re = 1000\).
\item
  Paroi supérieure mobile à \(u = 1.0\text{ m/s}\).
\end{itemize}

\subsubsection{5.2. Lancement}\label{lancement}

\begin{Shaded}
\begin{Highlighting}[]
\ExtensionTok{mpirun} \AttributeTok{{-}np}\NormalTok{ 4 bin/covof2d\_fortran cases/cavity\_re1000.json}
\end{Highlighting}
\end{Shaded}

\subsubsection{5.3. Résultats
Numériques}\label{ruxe9sultats-numuxe9riques}

Le vortex primaire central converge à la position
\((x, y) = (0.5898, 0.5977)\), avec capture précise des vortex
secondaires de coin sans oscillation numérique.

\begin{center}\rule{0.5\linewidth}{0.5pt}\end{center}

\subsection{6. Tutoriel 3 : Écoulement sur Marche Descendante (Armaly Re
=
100)}\label{tutoriel-3-uxe9coulement-sur-marche-descendante-armaly-re-100}

L'expérience d'Armaly et al.~(1983) étudie le décollement de la couche
limite laminaire et le recollement en aval d'une marche de hauteur
\(H = 1.0\text{ m}\).

\subsubsection{\texorpdfstring{6.1. Obstacle Immergé par Brinkman
(\texttt{cases/backward\_facing\_step.json})}{6.1. Obstacle Immergé par Brinkman (cases/backward\_facing\_step.json)}}\label{obstacle-immerguxe9-par-brinkman-casesbackward_facing_step.json}

La marche est définie dans le fichier JSON par pénalisation de Brinkman
:

\begin{Shaded}
\begin{Highlighting}[]
\StringTok{"obstacle"}\ErrorTok{:} \FunctionTok{\{}
  \DataTypeTok{"enabled"}\FunctionTok{:} \KeywordTok{true}\FunctionTok{,}
  \DataTypeTok{"type"}\FunctionTok{:} \StringTok{"BOX"}\FunctionTok{,}
  \DataTypeTok{"x0"}\FunctionTok{:} \FloatTok{0.0}\FunctionTok{,} \DataTypeTok{"x1"}\FunctionTok{:} \FloatTok{4.0}\FunctionTok{,}
  \DataTypeTok{"y0"}\FunctionTok{:} \FloatTok{0.0}\FunctionTok{,} \DataTypeTok{"y1"}\FunctionTok{:} \FloatTok{1.0}
\FunctionTok{\}}
\end{Highlighting}
\end{Shaded}

\subsubsection{6.2. Lancement}\label{lancement-1}

\begin{Shaded}
\begin{Highlighting}[]
\ExtensionTok{mpirun} \AttributeTok{{-}np}\NormalTok{ 4 bin/covof2d\_fortran cases/backward\_facing\_step.json}
\end{Highlighting}
\end{Shaded}

La longueur de recollement adimensionnelle obtenue est
\(x_r / H = 2.719\), en accord direct avec la valeur expérimentale de
référence (\(x_r / H = 2.808\), écart relatif de \(3.18\%\)).

\begin{center}\rule{0.5\linewidth}{0.5pt}\end{center}

\subsection{7. Tutoriel 4 : Convection Naturelle Thermique de Vahl Davis
(Ra =
10⁴)}\label{tutoriel-4-convection-naturelle-thermique-de-vahl-davis-ra-10ux2074}

Ce cas étudie la convection d'un fluide soumis à un gradient de
température horizontal entre deux parois isothermes (\(T_h = 1.0\),
\(T_c = 0.0\)), les frontières supérieure et inférieure étant
adiabatiques (\(\partial T / \partial y = 0\)).

\subsubsection{7.1. Modèle de Boussinesq}\label{moduxe8le-de-boussinesq}

Le couplage vitesse-température est modélisé par :
\[\mathbf{f}_{\text{boussinesq}} = -\rho \beta (T - T_0) \mathbf{g}\]

\subsubsection{7.2. Lancement et Nombre de
Nusselt}\label{lancement-et-nombre-de-nusselt}

\begin{Shaded}
\begin{Highlighting}[]
\ExtensionTok{mpirun} \AttributeTok{{-}np}\NormalTok{ 2 bin/covof2d\_fortran cases/natural\_convection\_ra1e4.json}
\end{Highlighting}
\end{Shaded}

Le nombre de Nusselt moyen mesuré à la paroi chaude
\(\overline{Nu} = 2.3990\) se situe parfaitement dans la plage de
référence Référence Académique CFD (\([2.2, 2.7]\)).

\begin{center}\rule{0.5\linewidth}{0.5pt}\end{center}

\subsection{8. Tutoriel 5 : Écoulement Diphasique VOF --- Rupture de
Barrage}\label{tutoriel-5-uxe9coulement-diphasique-vof-rupture-de-barrage}

L'effondrement gravitationnel d'une colonne d'eau (Martin \& Moyce) met
à l'épreuve l'algorithme VOF face aux forts gradients de vitesse et
déformations d'interface.

\begin{verbatim}
          t = 0.0 s                     t = 0.4 s (propagation front)
    +-----------------------+     +-----------------------+
    | ####                  |     |                       |
    | ####                  |     |                       |
    | ####                  |     | ##                    |
    | ####                  |     | #######======>        |
    +-----------------------+     +-----------------------+
\end{verbatim}

\subsubsection{\texorpdfstring{8.1. Configuration
(\texttt{cases/dam\_break.json})}{8.1. Configuration (cases/dam\_break.json)}}\label{configuration-casesdam_break.json}

\begin{itemize}
\tightlist
\item
  Fluide 1 (Eau) : \(\rho_1 = 1000\text{ kg/m}^3\),
  \(\mu_1 = 0.001\text{ Pa}\cdot\text{s}\).
\item
  Fluide 2 (Air) : \(\rho_2 = 1.0\text{ kg/m}^3\),
  \(\mu_2 = 1.85 \times 10^{-5}\text{ Pa}\cdot\text{s}\).
\item
  Schéma d'advection VOF : \textbf{TVD Superbee} garantissant la stricte
  conservation de volume et une interface de largeur maximale 2
  cellules.
\end{itemize}

\subsubsection{8.2. Lancement}\label{lancement-2}

\begin{Shaded}
\begin{Highlighting}[]
\ExtensionTok{mpirun} \AttributeTok{{-}np}\NormalTok{ 4 bin/covof2d\_fortran cases/dam\_break.json}
\end{Highlighting}
\end{Shaded}

\begin{center}\rule{0.5\linewidth}{0.5pt}\end{center}

\subsection{9. Tutoriel 6 : Tension de Surface et Saut de Pression de
Laplace}\label{tutoriel-6-tension-de-surface-et-saut-de-pression-de-laplace}

Ce cas valide le calcul de tension de surface par la méthode \textbf{CSF
(Continuum Surface Force)} pour une goutte statique de rayon
\(R = 0.25\text{ m}\).

\subsubsection{9.1. Saut de Pression
Théorique}\label{saut-de-pression-thuxe9orique}

\[\Delta P = P_{\text{int}} - P_{\text{ext}} = \frac{\sigma}{R}\] Pour
\(\sigma = 1.0\text{ N/m}\) et \(R = 0.25\text{ m}\), la théorie prédit
\(\Delta P = 4.000\text{ Pa}\).

\subsubsection{9.2. Lancement et
Validation}\label{lancement-et-validation}

\begin{Shaded}
\begin{Highlighting}[]
\ExtensionTok{mpirun} \AttributeTok{{-}np}\NormalTok{ 2 bin/covof2d\_fortran cases/static\_droplet.json}
\end{Highlighting}
\end{Shaded}

Le saut de pression moyen calculé par COVOF 2D Fortran est de
\(\Delta P = 3.86\text{ Pa}\) (écart relatif de \(3.38\%\)), confirmant
l'absence de courants parasites déstabilisants.

\begin{center}\rule{0.5\linewidth}{0.5pt}\end{center}

\subsection{10. Tutoriel 7 : Écoulement Autour d'un Cylindre par
Pénalisation de
Brinkman}\label{tutoriel-7-uxe9coulement-autour-dun-cylindre-par-puxe9nalisation-de-brinkman}

Ce tutoriel modélise l'écoulement stationnaire à \(Re=20\) autour d'un
cylindre de diamètre \(D=0.1\text{ m}\).

\subsubsection{\texorpdfstring{10.1. Configuration
(\texttt{cases/obstacle\_cylinder\_re20.json})}{10.1. Configuration (cases/obstacle\_cylinder\_re20.json)}}\label{configuration-casesobstacle_cylinder_re20.json}

\begin{Shaded}
\begin{Highlighting}[]
\StringTok{"obstacle"}\ErrorTok{:} \FunctionTok{\{}
  \DataTypeTok{"enabled"}\FunctionTok{:} \KeywordTok{true}\FunctionTok{,}
  \DataTypeTok{"type"}\FunctionTok{:} \StringTok{"CYLINDER"}\FunctionTok{,}
  \DataTypeTok{"xc"}\FunctionTok{:} \FloatTok{0.5}\FunctionTok{,} \DataTypeTok{"yc"}\FunctionTok{:} \FloatTok{0.5}\FunctionTok{,}
  \DataTypeTok{"radius"}\FunctionTok{:} \FloatTok{0.05}
\FunctionTok{\}}
\end{Highlighting}
\end{Shaded}

\subsubsection{10.2. Lancement}\label{lancement-3}

\begin{Shaded}
\begin{Highlighting}[]
\ExtensionTok{mpirun} \AttributeTok{{-}np}\NormalTok{ 2 bin/covof2d\_fortran cases/obstacle\_cylinder\_re20.json}
\end{Highlighting}
\end{Shaded}

Le sillage symétrique présente une longueur de recirculation
\(L/D = 0.768\), strictement conforme à la valeur de référence Référence
Académique CFD (\(0.768\)).

\begin{center}\rule{0.5\linewidth}{0.5pt}\end{center}

\subsection{11. Post-Traitement Avancé avec
ParaView}\label{post-traitement-avancuxe9-avec-paraview}

Le module \texttt{module\_vtk\_writer.f90} produit des sorties standards
XML VTK : * Fichier maître multi-processus :
\texttt{output\_cas/cas\_XXXXXX.pvtr} * Fichiers de sous-domaines
partitionnés : \texttt{output\_cas/cas\_p0\_XXXXXX.vtr},
\texttt{cas\_p1\_XXXXXX.vtr}, etc.

\subsubsection{Procédure de Visualisation
:}\label{procuxe9dure-de-visualisation}

\begin{enumerate}
\def\labelenumi{\arabic{enumi}.}
\tightlist
\item
  Lancer \textbf{ParaView} et ouvrir le groupe de fichiers
  \texttt{cas\_..pvtr}.
\item
  Cliquer sur \textbf{Apply}.
\item
  Sélectionner la variable \texttt{Velocity} et activer les vecteurs par
  le filtre \emph{Glyph}.
\item
  Tracer les lignes de courant en appliquant le filtre \emph{Stream
  Tracer}.
\item
  Isoler l'interface diphasique en appliquant le filtre \emph{Contour}
  sur le champ \texttt{ColorFunction} (\(C = 0.5\)).
\item
  Tracer le profil de vitesse 1D par le filtre \emph{Plot Over Line}.
\end{enumerate}

\begin{center}\rule{0.5\linewidth}{0.5pt}\end{center}

\subsection{12. Automatisation Python et Études de
Convergence}\label{automatisation-python-et-uxe9tudes-de-convergence}

Vous pouvez automatiser des calculs paramétriques et des études de
convergence de maillage via un script Python simple pilotant le binaire
Fortran :

\begin{Shaded}
\begin{Highlighting}[]
\ImportTok{import}\NormalTok{ subprocess}
\ImportTok{import}\NormalTok{ json}

\NormalTok{mesh\_sizes }\OperatorTok{=}\NormalTok{ [}\DecValTok{16}\NormalTok{, }\DecValTok{32}\NormalTok{, }\DecValTok{64}\NormalTok{, }\DecValTok{128}\NormalTok{]}
\NormalTok{results }\OperatorTok{=}\NormalTok{ []}

\ControlFlowTok{for}\NormalTok{ n }\KeywordTok{in}\NormalTok{ mesh\_sizes:}
    \ControlFlowTok{with} \BuiltInTok{open}\NormalTok{(}\StringTok{"cases/poiseuille\_variable\_viscosity.json"}\NormalTok{, }\StringTok{"r"}\NormalTok{) }\ImportTok{as}\NormalTok{ f:}
\NormalTok{        cfg }\OperatorTok{=}\NormalTok{ json.load(f)}
\NormalTok{    cfg[}\StringTok{"mesh"}\NormalTok{][}\StringTok{"ny"}\NormalTok{] }\OperatorTok{=}\NormalTok{ n}
\NormalTok{    cfg[}\StringTok{"output"}\NormalTok{][}\StringTok{"output\_dir"}\NormalTok{] }\OperatorTok{=} \SpecialStringTok{f"output\_conv\_}\SpecialCharTok{\{}\NormalTok{n}\SpecialCharTok{\}}\SpecialStringTok{"}
    
    \ControlFlowTok{with} \BuiltInTok{open}\NormalTok{(}\StringTok{"cases/tmp\_case.json"}\NormalTok{, }\StringTok{"w"}\NormalTok{) }\ImportTok{as}\NormalTok{ f:}
\NormalTok{        json.dump(cfg, f)}
        
\NormalTok{    subprocess.run([}\StringTok{"mpirun"}\NormalTok{, }\StringTok{"{-}np"}\NormalTok{, }\StringTok{"4"}\NormalTok{, }\StringTok{"bin/covof2d\_fortran"}\NormalTok{, }\StringTok{"cases/tmp\_case.json"}\NormalTok{])}
\end{Highlighting}
\end{Shaded}

\begin{center}\rule{0.5\linewidth}{0.5pt}\end{center}

\subsection{13. Créer un Nouveau Cas Test à Partir de
Zéro}\label{cruxe9er-un-nouveau-cas-test-uxe0-partir-de-zuxe9ro}

Pour simuler un écoulement inédit avec COVOF 2D Fortran :

\begin{enumerate}
\def\labelenumi{\arabic{enumi}.}
\item
  Créez un fichier \texttt{cases/mon\_cas.json}.
\item
  Définissez la grille : \texttt{nx}, \texttt{ny}, \texttt{lx},
  \texttt{ly}.
\item
  Précisez les propriétés physiques (\texttt{rho}, \texttt{mu},
  \texttt{gx}, \texttt{gy}).
\item
  Indiquez les conditions aux limites sur chaque paroi (\texttt{left},
  \texttt{right}, \texttt{bottom}, \texttt{top}).
\item
  Configurez le pas de temps (\texttt{dt}, \texttt{t\_final}) et les
  paramètres du solveur (\texttt{r\_al}, \texttt{linear\_solver}).
\item
  Exécutez le calcul :

\begin{Shaded}
\begin{Highlighting}[]
\ExtensionTok{mpirun} \AttributeTok{{-}np}\NormalTok{ 4 bin/covof2d\_fortran cases/mon\_cas.json}
\end{Highlighting}
\end{Shaded}
\item
  Visualisez immédiatement le résultat :

\begin{Shaded}
\begin{Highlighting}[]
\ExtensionTok{paraview}\NormalTok{ output\_mon\_cas/mon\_cas\_..pvtr}
\end{Highlighting}
\end{Shaded}
\end{enumerate}

\begin{center}\rule{0.5\linewidth}{0.5pt}\end{center}

\subsection{14. Tutoriel Avancé : Comparaison des Méthodes de Suivi
d'Interface
VOF}\label{tutoriel-avancuxe9-comparaison-des-muxe9thodes-de-suivi-dinterface-vof}

COVOF 2D intègre une architecture modulaire complète pour le suivi
d'interface diphasique, combinant plusieurs méthodes de reconstruction
géométrique d'interface et d'advection conservative : -
\textbf{Reconstruction d'interface} : \texttt{SLIC} (ordre 1
orthogonal), \texttt{PLIC} (ordre 2 linéaire continu, Youngs),
\texttt{LVIRA} (optimisation continue par moindres carrés),
\texttt{ELVIRA} (sélection optimale parmi fonctions de hauteur
discrètes). - \textbf{Algorithmes d'advection VOF} : \texttt{VOF\_TVD}
(algébrique compressif avec limiteurs), \texttt{CIAM} (géométrique par
découpage directionnel de Strang), \texttt{WEYMOUTH\_YUE} (géométrique
multidimensionnel non découplé avec correction de coin \(V_{xy}\)). -
\textbf{Limiteurs de flux TVD / NVD} : \texttt{SUPERBEE},
\texttt{VAN\_LEER}, \texttt{MINMOD}, \texttt{MUSCL}, \texttt{HRIC}.

\subsubsection{14.1. Le Banc d'Essai : Vortex Périodique Cisaillant de
LeVeque (Cas
3.19)}\label{le-banc-dessai-vortex-puxe9riodique-cisaillant-de-leveque-cas-3.19}

Dans ce cas test classique, une goutte circulaire de rayon
\(R_0 = 0.15\) initialement centrée en \((0.5, 0.75)\) dans un domaine
\([0, 1] \times [0, 1]\) est soumise à un champ de vitesse stationnaire
ou dépendant du temps à fort cisaillement :
\[u(x, y, t) = -2 \sin^2(\pi x) \sin(\pi y) \cos(\pi y) \cos\left(\frac{\pi t}{T}\right)\]
\[v(x, y, t) = 2 \sin(\pi x) \cos(\pi x) \sin^2(\pi y) \cos\left(\frac{\pi t}{T}\right)\]

La goutte s'étire en un filament spiralé extrêmement fin jusqu'à
\(t = T/2\), puis le champ de vitesse s'inverse pour ramener le fluide à
sa forme initiale circulaire à \(t = T\). La solution analytique à
\(t = T\) est rigoureusement identique à la condition initiale
\(C(\mathbf{x}, T) \equiv C(\mathbf{x}, 0)\).

L'erreur globale de réversibilité est mesurée par :
\[E_{L1} = \int_\Omega |C(\mathbf{x}, T) - C(\mathbf{x}, 0)| \, d\Omega = \sum_{i, j} |C_{i, j}^T - C_{i, j}^0| \Delta x \Delta y\]

\subsubsection{14.2. Configuration du Suivi d'Interface dans le Fichier
JSON}\label{configuration-du-suivi-dinterface-dans-le-fichier-json}

Pour sélectionner une combinaison spécifique dans
\texttt{cases/vortex\_shear.json}, modifiez la section \texttt{"vof"} :

\begin{Shaded}
\begin{Highlighting}[]
\StringTok{"vof"}\ErrorTok{:} \FunctionTok{\{}
  \DataTypeTok{"reconstruction\_method"}\FunctionTok{:} \StringTok{"ELVIRA"}\FunctionTok{,}
  \DataTypeTok{"advection\_method"}\FunctionTok{:} \StringTok{"WEYMOUTH\_YUE"}\FunctionTok{,}
  \DataTypeTok{"tvd\_limiter"}\FunctionTok{:} \StringTok{"SUPERBEE"}\FunctionTok{,}
  \DataTypeTok{"cfl\_vof"}\FunctionTok{:} \FloatTok{0.5}
\FunctionTok{\}}
\end{Highlighting}
\end{Shaded}

Options disponibles pour chaque paramètre : -
\texttt{"reconstruction\_method"} : \texttt{"SLIC"}, \texttt{"PLIC"},
\texttt{"LVIRA"}, \texttt{"ELVIRA"} - \texttt{"advection\_method"} :
\texttt{"VOF\_TVD"}, \texttt{"CIAM"}, \texttt{"WEYMOUTH\_YUE"} -
\texttt{"tvd\_limiter"} : \texttt{"SUPERBEE"}, \texttt{"VAN\_LEER"},
\texttt{"MINMOD"}, \texttt{"MUSCL"}, \texttt{"HRIC"}

\subsubsection{14.3. Exécution et Analyse Comparative
Pas-à-Pas}\label{exuxe9cution-et-analyse-comparative-pas-uxe0-pas}

Pour évaluer l'impact du schéma géométrique sur la préservation du
filament fluide :

\begin{enumerate}
\def\labelenumi{\arabic{enumi}.}
\item
  \textbf{Exécution avec le schéma algébrique VOF-TVD Superbee} :

\begin{Shaded}
\begin{Highlighting}[]
\CommentTok{\# Configurer: "reconstruction\_method": "PLIC", "advection\_method": "VOF\_TVD", "tvd\_limiter": "SUPERBEE"}
\ExtensionTok{mpirun} \AttributeTok{{-}np}\NormalTok{ 4 bin/covof2d\_fortran cases/vortex\_shear.json}
\end{Highlighting}
\end{Shaded}

  \emph{Résultat typique sur grille \(64 \times 64\) (\(T=4\text{ s}\))}
  : \(E_{L1} \approx 5.69 \times 10^{-2}\).
\item
  \textbf{Exécution avec le schéma géométrique CIAM + PLIC} :

\begin{Shaded}
\begin{Highlighting}[]
\CommentTok{\# Configurer: "reconstruction\_method": "PLIC", "advection\_method": "CIAM"}
\ExtensionTok{mpirun} \AttributeTok{{-}np}\NormalTok{ 4 bin/covof2d\_fortran cases/vortex\_shear.json}
\end{Highlighting}
\end{Shaded}

  \emph{Résultat typique sur grille \(64 \times 64\) (\(T=4\text{ s}\))}
  : \(E_{L1} \approx 6.23 \times 10^{-3}\) (\textbf{gain de précision
  d'un ordre de grandeur}).
\item
  \textbf{Exécution avec l'algorithme multidimensionnel de Weymouth \&
  Yue + ELVIRA} :

\begin{Shaded}
\begin{Highlighting}[]
\CommentTok{\# Configurer: "reconstruction\_method": "ELVIRA", "advection\_method": "WEYMOUTH\_YUE"}
\ExtensionTok{mpirun} \AttributeTok{{-}np}\NormalTok{ 4 bin/covof2d\_fortran cases/vortex\_shear.json}
\end{Highlighting}
\end{Shaded}

  \emph{Résultat typique sur grille \(64 \times 64\) (\(T=4\text{ s}\))}
  : \(E_{L1} \approx 1.18 \times 10^{-2}\) sans biais d'axe directionnel
  ni fragmentation du filament.
\end{enumerate}

\subsubsection{14.4. Tableau Synthétique des
Performances}\label{tableau-synthuxe9tique-des-performances}

{\def\LTcaptype{none} % do not increment counter
\begin{longtable}[]{@{}
  >{\raggedright\arraybackslash}p{(\linewidth - 10\tabcolsep) * \real{0.1947}}
  >{\raggedright\arraybackslash}p{(\linewidth - 10\tabcolsep) * \real{0.1416}}
  >{\raggedright\arraybackslash}p{(\linewidth - 10\tabcolsep) * \real{0.0885}}
  >{\raggedright\arraybackslash}p{(\linewidth - 10\tabcolsep) * \real{0.2655}}
  >{\raggedright\arraybackslash}p{(\linewidth - 10\tabcolsep) * \real{0.1770}}
  >{\raggedright\arraybackslash}p{(\linewidth - 10\tabcolsep) * \real{0.1327}}@{}}
\toprule\noalign{}
\begin{minipage}[b]{\linewidth}\raggedright
Algorithme Advection
\end{minipage} & \begin{minipage}[b]{\linewidth}\raggedright
Reconstruction
\end{minipage} & \begin{minipage}[b]{\linewidth}\raggedright
Limiteur
\end{minipage} & \begin{minipage}[b]{\linewidth}\raggedright
Erreur \(L_1\) Vortex (\(64^2\))
\end{minipage} & \begin{minipage}[b]{\linewidth}\raggedright
Conservation Masse
\end{minipage} & \begin{minipage}[b]{\linewidth}\raggedright
Temps relatif
\end{minipage} \\
\midrule\noalign{}
\endhead
\bottomrule\noalign{}
\endlastfoot
\texttt{VOF\_TVD} & N/A & \texttt{SUPERBEE} & \(5.69 \times 10^{-2}\) &
\(100.000\%\) & \(1.0\times\) (Référence) \\
\texttt{VOF\_TVD} & N/A & \texttt{VAN\_LEER} & \(7.12 \times 10^{-2}\) &
\(100.000\%\) & \(0.98\times\) \\
\texttt{VOF\_TVD} & N/A & \texttt{HRIC} & \(6.04 \times 10^{-2}\) &
\(100.000\%\) & \(1.05\times\) \\
\texttt{CIAM} & \texttt{SLIC} & N/A & \(2.41 \times 10^{-2}\) &
\(100.000\%\) & \(1.4\times\) \\
\texttt{CIAM} & \texttt{PLIC} & N/A & \(6.23 \times 10^{-3}\) &
\(100.000\%\) & \(1.8\times\) \\
\texttt{WEYMOUTH\_YUE} & \texttt{PLIC} & N/A & \(1.24 \times 10^{-2}\) &
\(100.000\%\) & \(2.1\times\) \\
\texttt{WEYMOUTH\_YUE} & \texttt{ELVIRA} & N/A & \(1.18 \times 10^{-2}\)
& \(100.000\%\) & \(2.6\times\) \\
\texttt{WEYMOUTH\_YUE} & \texttt{LVIRA} & N/A & \(1.15 \times 10^{-2}\)
& \(100.000\%\) & \(3.2\times\) \\
\end{longtable}
}

\subsubsection{14.5. Recommandations de Choix pour vos
Simulations}\label{recommandations-de-choix-pour-vos-simulations}

\begin{itemize}
\tightlist
\item
  \textbf{Grandes simulations industrielles / maillages très fins
  (\(> 512^2\))} : Utiliser \texttt{VOF\_TVD} avec limiteur
  \texttt{SUPERBEE} ou \texttt{HRIC}. Il offre le meilleur compromis
  coût CPU / robustesse et préserve rigoureusement la conservation de
  masse.
\item
  \textbf{Écoulements diphasiques fortement cisaillés ou microfluidique
  (bulles, gouttes)} : Utiliser \texttt{CIAM} avec \texttt{PLIC} ou
  \texttt{WEYMOUTH\_YUE} avec \texttt{ELVIRA} pour capturer la courbure
  et éliminer les courants parasites avec une précision optimale.
\end{itemize}

\end{document}
